% !TEX root = ../main.tex
% ESEMPIO 0b: slide Motivations con una figura animata (personaggi e bolle di dialogo).
%   La figura può essere un'immagine qualsiasi. Qui i numeri di sovrapposizione (overlay) sono definiti nel frame:
%   ogni passo parte dal precedente (+1) o resta uguale (+0) quando due cose compaiono insieme.
%   Gli stessi nomi sono usati da assets/tex/02_fig_cert_process.tex.
%   Da copiare: la struttura (macro dei numeri nel frame, figura in assets/tex/, \input dentro adjustbox).
%   La stessa tecnica si riusa nell'approccio, per mostrare una figura con i colori che evidenziano i passi.

\begin{frame}
    \frametitle{14 Legacy Nicola: animated scenario figure}
    % GUIDA: il problema, e perché lo stato dell'arte non basta. Meglio mostrarlo su un esempio concreto (una figura di uno scenario),
    % che può tornare più avanti nell'approccio.

    \pgfmathsetmacro\numberStart{int(1)}
    % passo: si evidenzia l'elemento principale e le sue dipendenze
    \pgfmathsetmacro\numberServiceTopStart{int(\numberStart+0)}
    \pgfmathsetmacro\numberHighlightDependencyStart{int(\numberServiceTopStart+0)}
    % passo: nuova bolla sull'elemento principale, che poi smette di essere evidenziato
    \pgfmathsetmacro\numberServiceTopEvidStart{int(\numberHighlightDependencyStart+1)}
    \pgfmathsetmacro\numberServiceTopEvidEnd{int(\numberServiceTopEvidStart+0)}
    \pgfmathsetmacro\numberServiceTopEnd{int(\numberServiceTopEvidEnd+0)}
    % passo: restano evidenziati gli altri elementi, nuova bolla
    \pgfmathsetmacro\numberCalloutFocusOnOtherStart{int(\numberServiceTopEvidEnd+1)}
    \pgfmathsetmacro\numberCalloutFocusOnOtherEnd{int(\numberCalloutFocusOnOtherStart+0)}
    % passo: doppio evidenziamento di più elementi
    \pgfmathsetmacro\numberHighlightDependencyEnd{int(\numberCalloutFocusOnOtherEnd+0)}
    \pgfmathsetmacro\numberHighlightDependencySomeOkStart{int(\numberCalloutFocusOnOtherStart+0)}
    \pgfmathsetmacro\numberHighlightDependencySomeOkEnd{int(\numberHighlightDependencySomeOkStart+0)}
    % passo: doppio evidenziamento di un altro elemento, nuova bolla
    \pgfmathsetmacro\numberHighlightDependencySomeNotOkStart{int(\numberHighlightDependencySomeOkEnd+1)}
    \pgfmathsetmacro\numberHighlightDependencySomeNotOkStartPartOne{int(\numberHighlightDependencySomeNotOkStart+0)}
    \pgfmathsetmacro\numberHighlightDependencySomeNotOkStartPartOneEnd{int(\numberHighlightDependencySomeNotOkStartPartOne+0)}
    \pgfmathsetmacro\numberHighlightDependencySomeNotOkStartPartTwoStart{int(\numberHighlightDependencySomeNotOkStartPartOneEnd+0)}
    \pgfmathsetmacro\numberHighlightDependencySomeNotOkStartPartTwoEnd{int(\numberHighlightDependencySomeNotOkStartPartTwoStart+0)}
    \pgfmathsetmacro\numberHighlightDependencySomeNotOkEnd{int(\numberHighlightDependencySomeNotOkStartPartTwoEnd+0)}
    \pgfmathsetmacro\numberEnd{int(\numberHighlightDependencySomeNotOkEnd+0)}
    % passo finale: ogni elemento mostra la sua bolla e compare l'immagine
    \pgfmathsetmacro\numberCertStart{int(\numberHighlightDependencySomeNotOkEnd+1)}

    \pgfmathsetmacro\numberPatrickHappy{int(\numberCertStart+1)}
    \pgfmathsetmacro\numberCertEnd{int(\numberPatrickHappy+0)}

    \begin{adjustbox}{max totalsize={.9\linewidth}{\textheight},center}
        
	
\begin{tikzpicture}
			% \tikzset{
			% 	every node/.style={font=\scriptsize}
			% }
	\node[service in certification process,
    on slide={<\numberServiceTopStart> highlight light},
    on slide={<\numberServiceTopEvidStart-\numberServiceTopEvidEnd> highlight light},
    ] (target) {\makecell{Data Analysis Engine (\target)}};
	
    \node[service,below=0.6 of target,
    on slide={<\numberServiceTopStart> highlight light},
    on slide={<\numberCalloutFocusOnOtherStart-\numberCalloutFocusOnOtherEnd> highlight light}
    ] (auth) {Authentication};

	\node[service,below=0.5 of auth,
    on slide={<\numberServiceTopStart> highlight light},
    on slide={<\numberCalloutFocusOnOtherStart-\numberCalloutFocusOnOtherEnd> highlight light}
    ] (db) {Database};

	\node[service,below=0.5 of db,
    on slide={<\numberServiceTopStart> highlight light},
    on slide={<\numberCalloutFocusOnOtherStart-\numberCalloutFocusOnOtherEnd> highlight light}
    ] (ds) {\makecell{Distributed\\Storage}};
	
    \node[service,left=1.5 of ds,
    on slide={<\numberServiceTopStart> highlight light},
    on slide={<\numberCalloutFocusOnOtherStart-\numberCalloutFocusOnOtherEnd> highlight light}
    ] (mn) {\makecell{Mesh\\Network}};
	
    \node[service,right=2 of ds,
    on slide={<\numberServiceTopStart> highlight light},
    on slide={<\numberHighlightDependencySomeNotOkStart-\numberHighlightDependencySomeNotOkEnd> highlight light}
    ] (k8s) {Kubernetes};
	
    \node[service,below=0.7 of mn,
    on slide={<\numberServiceTopStart> highlight light},
    on slide={<\numberCalloutFocusOnOtherStart-\numberCalloutFocusOnOtherEnd> highlight light}
    ] (ni) {\makecell{Network Interface}};
	
    \node[service,below=0.7 of ds,
    on slide={<\numberServiceTopStart> highlight light},
    on slide={<\numberCalloutFocusOnOtherStart-\numberCalloutFocusOnOtherEnd> highlight light}
    ] (si) {\makecell{Storage Interface}};

	\begin{scope}[every node/.style={font=\footnotesize}]
		\draw[->, link] (target) -- node [left] {\tiny depends on} (auth);
		\draw[->, link] ([yshift=-5pt]target.east) -- ++(.25,0) |- node [left,pos=.175] {\tiny depends on} ([yshift=3pt]db.east);
		\draw[->, link] (auth) -- node [left] {\tiny depends on} (db);
		\draw[->, link] (db) -- node [left] {\tiny depends on} (ds);
		\draw[->, link] (ds) -- node [above] {\tiny depends on} (mn);
		\draw[->, link] (target) -| node [above,pos=.25] {\tiny depends on} ([xshift=-5pt]mn.north);
		\draw[->, link] (auth) -| node [above,pos=.25] {\tiny depends on} (mn.north);
		\draw[->, link] (db) -| node [above,pos=.25] {\tiny depends on} ([xshift=5pt]mn.north);
		\draw[->, link] ([yshift=-3pt]k8s.west) -- node [below] {\tiny coordinates} ([yshift=-3pt]ds.east);
		\draw[->, link] ([xshift=-5]k8s.north) |- node [below left] {\tiny coordinates} ([yshift=-3pt]db.east);
		\draw[->, link] (k8s.north) |- node [below left] {\tiny coordinates} (auth.east);
		\draw[->, link] ([xshift=5]k8s.north) |- node [below left] {\tiny coordinates} ([yshift=5pt]target.east);
		\draw[->, link] ([xshift=3]k8s.south) -- ++(0,-.6) -| node [below,pos=.25] {\tiny coordinates} ([xshift=3pt]si.north);
		\draw[->, link] ([xshift=-3]k8s.south) -- ++(0,-.4) -| node [below right] {\tiny coordinates} ([xshift=3pt]ni.north);
		\draw[->, link] (target.east) -- ++(.5,0) |- node [right,pos=.128] {\tiny depends on} ([yshift=3pt]ds.east);
		\draw[->, link] ([xshift=-3]ds.south) -- node [left] {\tiny depends on} ([xshift=-3pt]si.north);
		\draw[->, link] ([xshift=-3]mn.south) -- node [left] {\tiny depends on} ([xshift=-3pt]ni.north);
	\end{scope}


    \path let \p1 = (target), \p2 = (k8s) in
        coordinate (person-1-pos) at (\x2, \y1);
    \coordinate (person-1-pos) at ([xshift=20]person-1-pos); 

    \node[dave, minimum size=.75cm, xshift=20,
        visible on=<\numberStart-\numberEnd>
    ] (person-1) at (person-1-pos) {};
    %
    %
    % 2. let's certify it.
    \node[speaking question, callout absolute pointer=(person-1.mouth), anchor=west,
    visible on=<\numberServiceTopStart>
    ] (person-1-speak-1) at ([xshift=15]person-1.east) {%
        \begin{varwidth}{\widthof{These are its dependencies, }}
            \raggedright
            Let's certify this service. These are its dependencies, specified in \certModel.
        \end{varwidth}%
    };
    %
    %
    % 3. collect evidence from there
    \node[speaking question, callout absolute pointer=(person-1.mouth), anchor=west,
    visible on=<\numberServiceTopEvidStart-\numberServiceTopEvidEnd>
    ] (person-1-speak-2) at ([xshift=15]person-1.east) {%
        \begin{varwidth}{\widthof{These are its dependencies, }}
            \raggedright
            Let's certify this service. First, we collect evidence from its mechanisms.
        \end{varwidth}%
    };
    %
    %
    % 4. exists cert compatible?
    \node[speaking question, callout absolute pointer=(person-1.mouth), anchor=west,
    visible on=<\numberCalloutFocusOnOtherStart-\numberCalloutFocusOnOtherEnd>
    ] (person-1-speak-3) at ([xshift=15]person-1.east) {%
        \begin{varwidth}{\widthof{These are its dependencies, }}
            \raggedright
            Let's certify this service. \alert{These are its dependencies: are they already certified according to $\certModel.\referred{\certModel_i}$}?
        \end{varwidth}%
    };
    %
    %
    % 5. highlight many that are
    % so we need a judge
    \node[judge, minimum size=.75cm, xshift=20,
        %visible on=<\numberHighlightDependencySomeOkStart-\numberHighlightDependencySomeOkEnd>
        visible on=<\numberCalloutFocusOnOtherStart-\numberCalloutFocusOnOtherEnd>
    ] (person-2) at ([yshift=-80]person-1-pos) {};
    \node[speaking question, callout absolute pointer=(person-2.mouth), anchor=west,
    %visible on=<\numberHighlightDependencySomeOkStart-\numberHighlightDependencySomeOkEnd>
        visible on=<\numberCalloutFocusOnOtherStart-\numberCalloutFocusOnOtherEnd>
    ] (person-2-speak-4) at ([xshift=15]person-2.east) {%
        \begin{varwidth}{\widthof{I hereby guarantee that }}
            \raggedright
            I hereby guarantee that these dependencies are already certified.
        \end{varwidth}%
    };
    %
    %
    % 6. pt1 highlight some that are not.
    % we need a judge
    \node[judge, minimum size=.75cm, xshift=20,
    visible on=<\numberHighlightDependencySomeNotOkStart-\numberHighlightDependencySomeNotOkEnd>
    ] (person-5) at (person-1-pos) {};
    \node[speaking question, callout absolute pointer=(person-1.mouth), anchor=west,
    visible on=<\numberHighlightDependencySomeNotOkStart-\numberHighlightDependencySomeNotOkEnd>
    ] (person-2-speak-5) at ([xshift=15]person-1.east) {%
        \begin{varwidth}{\widthof{I hereby declare that}}
            \raggedright
            I hereby declare that these dependencies are not already certified.
        \end{varwidth}%
    };
    % we need dave right below
    \node[dave, minimum size=.75cm, xshift=20,
    visible on=<\numberHighlightDependencySomeNotOkStart-\numberHighlightDependencySomeNotOkEnd>
    ] (person-3) at ([yshift=-80]person-1-pos) {};
    %
    \node[speaking question, callout absolute pointer=(person-3.mouth), anchor=west,
    visible on=<\numberHighlightDependencySomeNotOkStart-\numberHighlightDependencySomeNotOkEnd>
    ] (person-3-speak-5) at ([xshift=15]person-3.east) {%
        \begin{varwidth}{\widthof{Let's repeat}}
            \raggedright
            Let's repeat the process recursively.
        \end{varwidth}%
    };
    %
    %
    % 7. certificate
    % we re-add the judge just because
    \node[judge, minimum size=.75cm, xshift=20,
        visible on=<\numberCertStart-\numberCertEnd>
    ] (person-4) at ([yshift=-50]person-1-pos) {};
    %
    \node[speaking question, callout absolute pointer=(person-4.mouth), anchor=west,
    visible on=<\numberCertStart-\numberCertEnd>
    ] (person-3-speak-5) at ([xshift=15]person-4.east) {%
        \begin{varwidth}{\widthof{Based on the evidence and the }}
            \raggedright
            Based on the evidence and the referred certificates provided to me, I hereby award a \alert{certificate \cert}, including the evidence, the \alert{referred certificates $\{\referred{\cert}\}$}, and my signature. 
        \end{varwidth}%
    };


    % south
    \foreach \src/\pointer in {ni/ni.south, si/si.south}{
        \node[speaking question, callout absolute pointer=(\pointer), anchor=north,
        visible on=<\numberCertStart-\numberCertEnd>,
        %visible on=<\numberCalloutFocusOnOtherStart-\numberCalloutFocusOnOtherEnd>
        ] at ([yshift=-10]\pointer) (\src-cert) {%
            %\begin{varwidth}{\widthof{certified}}
                \tiny I am certified
            %\end{varwidth}
        };
        \node[speaking question, callout absolute pointer=(\pointer), anchor=north,
        %visible on=<\numberCertStart-\numberCertEnd>,
        visible on=<\numberCalloutFocusOnOtherStart-\numberCalloutFocusOnOtherEnd>
        ] at ([yshift=-10]\pointer) (\src-cert-2) {%
            %\begin{varwidth}{\widthof{certified}}
                \tiny I am certified
            %\end{varwidth}
        };
    }
    % mesh
    \node[speaking question, callout absolute pointer=(mn.west), anchor=east,
        visible on=<\numberCertStart-\numberCertEnd>,
        %visible on=<\numberCalloutFocusOnOtherStart-\numberCalloutFocusOnOtherEnd>
        ] (mn-cert) at ([xshift=-10]mn.west) {%
            %\begin{varwidth}{\widthof{certified}}
                \tiny I am certified
            %\end{varwidth}
        };
    \node[speaking question, callout absolute pointer=(mn.west), anchor=east,
        %visible on=<\numberCertStart-\numberCertEnd>,
        visible on=<\numberCalloutFocusOnOtherStart-\numberCalloutFocusOnOtherEnd>
        ] (mn-cert) at ([xshift=-10]mn.west) {%
            %\begin{varwidth}{\widthof{certified}}
                \tiny I am certified
            %\end{varwidth}
        };
    % db
    \path let \p1 = (mn-cert), \p2 = (db) in
        node[speaking question, callout absolute pointer=(db.west), %anchor=east,
         yshift=5,
        visible on=<\numberCertStart-\numberCertEnd>,
        %visible on=<\numberCalloutFocusOnOtherStart-\numberCalloutFocusOnOtherEnd>
        ] (db-cert) at (\x1, \y2) {%
            \tiny I am certified
        };
    % repeat.
    \path let \p1 = (mn-cert), \p2 = (db) in
        node[speaking question, callout absolute pointer=(db.west), %anchor=east,
         yshift=5,
        %visible on=<\numberCertStart-\numberCertEnd>,
        visible on=<\numberCalloutFocusOnOtherStart-\numberCalloutFocusOnOtherEnd>
        ] (db-cert-2) at (\x1, \y2) {%
            \tiny I am certified
        };
    % auth
    \path let \p1 = (db-cert), \p2 = (auth) in
        node[speaking question, callout absolute pointer=(auth.west), %anchor=east,
         yshift=5,
        %visible on=<\numberCalloutFocusOnOtherStart-\numberCalloutFocusOnOtherEnd>
        visible on=<\numberCertStart-\numberCertEnd>
        ] at (\x1, \y2) (auth-cert) {%
            \tiny I am certified
        };
    \path let \p1 = (db-cert), \p2 = (auth) in
        node[speaking question, callout absolute pointer=(auth.west), %anchor=east,
         yshift=5,
        visible on=<\numberCalloutFocusOnOtherStart-\numberCalloutFocusOnOtherEnd>
        %visible on=<\numberCertStart-\numberCertEnd>
        ] at (\x1, \y2) (auth-cert-2) {%
            \tiny I am certified
        };
    % target
    \node[speaking question, callout absolute pointer=(target.north), anchor=south,
        visible on=<\numberCertStart-\numberCertEnd>
        ] at ([yshift=10]target.north) (target-cert) {%
            \tiny I am certified
        };
    % ds
    \node[speaking question, callout absolute pointer=(ds.south east), anchor=north west,
        visible on=<\numberCertStart-\numberCertEnd>,
        %visible on=<\numberCalloutFocusOnOtherStart-\numberCalloutFocusOnOtherEnd>
        ] at (si -| k8s) (ds-cert) {%
            \tiny I am certified
        };
    \node[speaking question, callout absolute pointer=(ds.south east), anchor=north west,
        %visible on=<\numberCertStart-\numberCertEnd>,
        visible on=<\numberCalloutFocusOnOtherStart-\numberCalloutFocusOnOtherEnd>
        ] at (si -| k8s) (ds-cert-2) {%
            \tiny I am certified
        };  
    % k8s
    \node[speaking question, callout absolute pointer=(k8s.east), anchor=west,
        visible on=<\numberCertStart-\numberCertEnd>
        ] at ([xshift=10]k8s.east) (k8s-cert) {%
            \tiny I am certified
        }; 


    \node[visible on=<\numberPatrickHappy>, anchor=north] at ([yshift=-10]person-3-speak-5.south) {\includegraphics[width=.3\textwidth]{assets/img/patrick_happy.jpg}};


\end{tikzpicture}


    \end{adjustbox}

\end{frame}
